% Direct-authority diplomatic French for printed p.204.
\clearpage
\SourcePage{204}

\LayerTarget{tohoku:remark:condition-D-basic-properties:p1}
\noindent topologie de Zariski et $G$ un groupe fini d'automorphismes de $X$,
la condition (D) n'est pas vérifiée (sauf si $G$ opère trivialement).

\LayerTarget{tohoku:theorem:5.3.1}
\noindent\textsc{Théorème 5.3.1.} \textit{Supposons la condition (D) vérifiée.
Soit $A$ un $G$-faisceau abélien sur $X$, soit $y\in Y$, $x$ un élément de
$f^{-1}(y)$. On a alors des isomorphismes canoniques}

\LayerTarget{tohoku:equation:5.3.1}
\[
\mathbf H^n(G,A)(y)=H^n(G_x,A(x)).
\tag{5.3.1}
\]

\LayerTarget{tohoku:proof:theorem-5.3.1}
Prenant $V_x$ comme dans l'énoncé de la condition (D), et de plus $V_x$
ouvert et stable par $G_x$, on peut appliquer le corollaire à la proposition
5.2.1, ce qui nous ramène au cas où $G=G_x$, $X=V_x$ donc où $G$ est
\textit{fini}. Montrons qu'alors $R^qf_*(A)=0$ pour $q>0$. On a pour tout
$y\in Y$ :

\LayerTarget{tohoku:equation:5.3.1-stalk-limit}
\[
R^qf_*(A)(y)=\varinjlim H^q(f^{-1}(U),A),
\]
la limite étant prise suivant les voisinages ouverts $U$ de $y$. Or $x$ et
$V_x$ étant pris comme ci-dessus, et se bornant aux $U\subset f(V_x)$, on
voit que $f^{-1}(U)$ est réunion finie d'ouverts disjoints deux à deux
$g\cdot U_x$ ($g\in G/G_x$), où $U_x=V_x\cap f^{-1}(U)$, donc
$H^q(f^{-1}(U),A)$ est somme directe d'un nombre fini de groupes isomorphes à
$H^q(U_x,A)$. Quand $U$ parcourt un système fondamental de voisinages de $y$,
$U_x$ parcourt un système fondamental de voisinages de $x$, donc on obtient
pour $R^qf_*(A)$ la somme directe d'un nombre fini de groupes isomorphes à
$\varinjlim H^q(U_x,A)$ ; or ce dernier est nul pour $q>0$, en vertu du lemme
3.8.2.

On a donc bien $R^qf_*(A)=0$ pour $q>0$, donc en vertu de proposition 5.2.2 le
premier membre de (5.3.1) s'identifie à $\mathbf H^n(G,f_*(A))(y)$, qui, $G$
étant fini, est égal à $H^n(G,f_*(A)(y))$ (comme on l'a signalé dans 4.4). Se
rappelant qu'on pouvait supposer $X=V_x$, donc que $f^{-1}(y)$ est réduit à
$x$, on a alors $f_*(A)(y)=A(x)$, ce qui achève la démonstration.

\LayerTarget{tohoku:corollary:5.3.1:1}
\noindent\textsc{Corollaire.} \textit{Supposons que pour tout $x\in X$,
l'ordre $n_x$ du stabilisateur $G_x$ soit tel que la multiplication par $n_x$
soit un automorphisme du groupe $A(x)$. Alors
$H^*(X;G,A)=H^*(Y,A^G)$, et par suite $H^*(Y,A^G)$ est l'aboutissement d'une
suite spectrale dont le terme initial est
$\mathrm{II}_2^{p,q}=H^p(G,H^q(X,A))$.}

\LayerTarget{tohoku:proof:corollary-5.3.1:1}
En effet, le théorème 3.5.1 donne $\mathbf H^q(G,A)=0$ pour $q>0$, donc on est
dans les conditions d'application de la proposition 5.2.3. Le corollaire
précédent est surtout intéressant dans le cas où $A$ est un faisceau d'espaces
vectoriels sur un corps de caractéristique $p$ ne divisant aucun des $n_x$
(ce qui est toujours le cas en caractéristique 0). Quand il n'y a pas de corps
de base, on pourra utiliser la variante suivante :

\LayerTarget{tohoku:corollary:5.3.1:2}
\noindent\textsc{Corollaire 2.} \textit{Supposons que les ordres $n_x$ des
$G_x$ admettent un ppcm $n$ ; soit $P$ l'ensemble des diviseurs premiers de
$n$, $C(P)$ la catégorie des groupes abéliens dont chaque élément a un ordre
fini n'ayant pas d'autres diviseurs premiers que des $p\in P$. Alors avec la
terminologie introduite dans 1.11, $H^*(X;G,A)$ est isomorphe mod $C(P)$ à
$H^*(Y,A^G)$, donc ce dernier est l'aboutissement mod $C(P)$ d'une suite
spectrale cohomologique dont le terme initial est $H^p(G,H^q(X,A))$.}

\LayerTarget{tohoku:proof:corollary-5.3.1:2:p1}
En effet, réduisons mod $C(P)$ les suites spectrales du théorème 5.2.1.

\TohokuContinuation{205}
